\documentclass[10pt,a4paper]{amsart}
\usepackage[margin=19mm]{geometry}
\usepackage{amsmath,amssymb,mathtools}
\usepackage[scr=rsfs,cal=euler]{mathalfa}
\usepackage{xfrac}
\usepackage[unicode,hidelinks,bookmarks=false]{hyperref}
\newcommand{\ie}{i.e.\ }
\newcommand{\cf}{cf.\ }
\newcommand{\resp}{respectively\ }
\newcommand{\Subsection}{\subsection}
\providecommand{\nobreakdash}{\nobreak\hbox{-}}
\setlength{\emergencystretch}{2em}
\multlinegap0pt
\usepackage{bm}
\usepackage{bbm}
\usepackage{dsfont}
\usepackage{xspace}
\usepackage{cancel}
\usepackage{mathtools}
\usepackage{amsthm}
\makeatletter
\@ifundefined{theorem}{\newtheorem{theorem}{Theorem}}{}
\@ifundefined{lemma}{\newtheorem{lemma}[theorem]{Lemma}}{}
\@ifundefined{proposition}{\newtheorem{proposition}[theorem]{Proposition}}{}
\makeatother
\DeclareMathOperator{\sgn}{sgn}
\newcommand{\Z}{\mathbb{Z}}
\newcommand{\cT}{\mathcal{T}}
\title[Construction of Generically Ordinary Families of Hyperellipt]{Construction of Generically Ordinary Families of Hyperelliptic Curves}
\author{Hui June Zhu}
\date{Source: arXiv:2606.31783v3 (CC BY 4.0)}
\begin{document}
\maketitle

\noindent\emph{Source and licence.} Construction of Generically Ordinary Families of Hyperelliptic Curves. Version arXiv:2606.31783v3 (CC BY 4.0), distributed under the Creative Commons Attribution 4.0 International licence, as recorded in the arXiv licence metadata for this exact version and shown on the abstract page https://arxiv.org/abs/2606.31783v3. Full rights notice: licenses/arxiv-2606.31783-CC-BY-4.0.txt.\par\smallskip

\noindent\textit{Editorial scope.} Counted unit: the Proposition whose source label is \texttt{P:D\_p} in the article (source0.tex lines 1004--1022, source label \texttt{P:D\_p}), delivered together with the article's own complete original proof of it (source0.tex lines 1024--1212, one \texttt{proof} environment of 189 lines). No mathematics is written, repaired or re-derived: statement and proof are the article's own text line for line. Required context retained uncounted: E:eta (lines 819--824); L:entry (lines 872--886); L:E (lines 974--980). Every definition, display and label of those lines is retained unchanged. Omitted from this package and not redistributed: 1--818, 825--871, 887--973, 981--1003, 1023--1023, 1213--1539. Nothing inside a retained excerpt is omitted or altered: every retained body is delivered exactly as the article's lines are written, so no omission receipt of any kind accompanies this package. Citations: the retained excerpts carry no citation command at all, so no bibliography entry of the article is transcribed and no reference is redistributed. Reference adaptation: none was needed and none was made; every reference inside the delivered unit resolves inside it, so no unresolved reference or citation is produced. Printed slips kept literal: no printed word, symbol, hypothesis or number of the counted statement or of its proof was altered. Layout and class: the article's own class is \texttt{amsart} with packages including amsmath, hyperref, amsfonts, amssymb, amsthm, mathrsfs, mathtools, comment, pdfpages, inputenc, graphicx, fontenc, lmodern; the wrapper is the lane's pinned \texttt{amsart} editorial wrapper at 10pt on A4, which restates the theorem-like environments the retained text needs and adds the article's own macro definitions that the retained text needs (\texttt{\textbackslash Z}, \texttt{\textbackslash cT}, \texttt{\textbackslash sgn}), with extra wrapper packages (bm, bbm, dsfont, xspace, cancel, mathtools). The delivered numbering is the unit's own. Assets: no retained span contains a figure environment, a graphics-inclusion command, a PDF-page inclusion, a table or a TikZ picture; no image file is copied into this package, the package declares no asset dependency and no image file is redistributed with it. Licence: version arXiv:2606.31783v3 of this article is distributed under the Creative Commons Attribution 4.0 International licence, as recorded in the arXiv licence metadata for this exact version and shown on the abstract page https://arxiv.org/abs/2606.31783v3. A full rights notice is delivered at licenses/arxiv-2606.31783-CC-BY-4.0.txt. Characters: every retained excerpt is pure ASCII, so the delivered engine, XeLaTeX with its default Latin Modern text font, reports no missing character.
\begin{equation}\label{E:eta}
q_{ij}=Q_i-\delta_{ij},
\qquad
r_{ij}=s_i-j+d\delta_{ij},\qquad
\eta_{ij}=\frac{p-1}{2}-q_{ij}-r_{ij}.
\end{equation}

\begin{lemma}\label{L:entry}
If $\eta_{ij}<0$ then $M_{ij}(X,Y)=0$. 
We have $\cT(p\,i-j)\ne \varnothing$ if and only if $\eta_{ij}\ge 0$. 
If $\eta_{ij}\ge 0$, then 
$$
\deg_Y(M_{ij}(X,Y))=  \eta_{ij},\qquad
[Y^{\eta_{ij}}]M_{ij}(X,Y)=X^{r_{ij}}b_{ij}.
$$
In particular, for every algebraic integer $\alpha$, 
$$\deg_Y M_{ij}(\alpha,Y)\le \eta_{ij},
\qquad
[Y^{\eta_{ij}}]M_{ij}(\alpha,Y)=\alpha^{r_{ij}}b_{ij};
$$
if $\alpha\ne 0$, then equality holds in the degree bound.
\end{lemma}

\begin{lemma}\label{L:E}
We have $S_g'\ne \varnothing$, and $E(\sigma)$ is independent of 
$\sigma\in S_g'$. Then 
$
E=\sum_{i=1}^g \left(\frac{p-1}{2}-Q_i-s_i\right)+\frac{g(g+1)}{2}.
$
\end{lemma}

\begin{proposition}\label{P:D_p} %Proposition 3.7
Let
$m=\sum_{i=1}^gs_i-\frac{g(g+1)}{2}$ and $u_p=
    \prod_{i=1}^g\binom{\frac{p-1}{2}}{Q_i}$.
Then $m\ge 0$ and $u_p$ is a positive integer coprime to $p$. 
\begin{enumerate}
\item If $E\ge 0$ then 
$$\deg_Y(H_{p}(X,Y))\le E, \qquad
[Y^E]H_{p}(X,Y) =X^m u_p \Delta_p.$$
\item Let $E\ge 0$. For every commutative ring $R$ and every $a \in R$, the specialization $H_p(a,Y)\in R[Y]$ satisfies 
$$
\deg_YH_p(a,Y)\le E, \qquad 
[Y^E]H_p(a,Y)=a^m u_p \Delta_p.
$$
In particular, if $m=0$, we have $[Y^E]H_p(0,Y)=u_p\Delta_p.$
\item 
If $E<0$ then $H_{p}(X,Y)=0$ and $\Delta_p=0$.
\end{enumerate}
\end{proposition}

\begin{proof}
Since $s_1,\ldots,s_g$ are distinct positive integers, 
$
\sum_{i=1}^{g}s_i\geq \sum_{i=1}^{g}i=\frac{g(g+1)}2,
$
and hence $m\geq0$.
Set
$
n=\frac{p-1}{2}.
$
Since $1\leq i\leq g<d/2$, we have
$
Q_i=\left\lfloor\frac{pi}{d}\right\rfloor<\frac p2.
$
Because $Q_i$ is an integer, it follows that
$
0\leq Q_i\leq \frac{p-1}{2}=n.
$
Thus $\binom{n}{Q_i}$ is a positive integer. Moreover, $n<p$, so none of the factorials occurring in
$
\binom{n}{Q_i}=\frac{n!}{Q_i!(n-Q_i)!}
$
is divisible by $p$. Hence each 
$\binom{n}{Q_i}$, and therefore their product $u_p$, is coprime to $p$.

\noindent
{\bf(1).} 
By Lemma \ref{L:entry}, if $\eta_{ij}<0$, then  $M_{ij}(X,Y)=0$.
If $\eta_{ij}\ge 0$, then the unique element of $\cT(pi-j)$ with maximal
$t_0$-coordinate is 
$(t_0,t_1,t_d)=(\eta_{ij},r_{ij},q_{ij})$.
For a permutation $\sigma\in S_g$, define
$
e(\sigma):=\#\{i:\sigma(i)>s_i\}.
$
Using the definition of $\eta_{ij}$ in \eqref{E:eta} and the equality
$
\sum_{i=1}^{g}\sigma(i)=\frac{g(g+1)}2,
$
we obtain
\begin{align}\label{E:eta_sum}
\sum_{i=1}^{g}\eta_{i,\sigma(i)}
&=
\sum_{i=1}^{g}
\left(
n-q_{i,\sigma(i)}-r_{i,\sigma(i)}\right)=
E-(d-1)e(\sigma).
\end{align}
Suppose first that $E\ge 0$. Consider the Leibniz expansion
$$
H_p(X,Y)=\sum_{\sigma\in S_g}\sgn(\sigma)\prod_{i=1}^g M_{i,\sigma(i)}(X,Y).
$$
Whenever the product corresponding to $\sigma$ is nonzero, all its
factors are nonzero, and hence by \eqref{E:eta_sum},
$$
\deg_Y\prod_{i=1}^g M_{i,\sigma(i)}(X,Y)=
\sum_{i=1}^{g}\eta_{i,\sigma(i)}
=E-(d-1)e(\sigma)\leq E.
$$
Thus $\deg_YH_p(X,Y)\le E$.
Moreover, if $e(\sigma)>0$, then every nonzero product corresponding to $\sigma$ has degree $<E$. Therefore, only permutations in 
$$
S'_g=\{\sigma\in S_g: \sigma(i)\leq s_i\text{ for every $i$}\}
$$
can contribute to $[Y^E]H_p(X,Y)$.
Let $\sigma\in S_g'$. 
If $\prod_{i=1}^g M_{i,\sigma(i)}(X,Y)\ne 0$, 
then 
\begin{align}
[Y^E]\prod_{i=1}^{g} M_{i,\sigma(i)}(X,Y)
&=X^{\sum_{i=1}^g (s_i-\sigma(i))} \prod_{i=1}^gb_{i,\sigma(i)}
=X^m \prod_{i=1}^g b_{i,\sigma(i)}.
\label{E:11}
\end{align}
This formula remains valid if one of the factors 
$M_{i,\sigma(i)}(X,Y)
$
is zero. Indeed, in that case $\eta_{i,\sigma(i)}<0$ for some $i$.
Since $\sigma(i)\le s_i$,
$
\eta_{i,\sigma(i)}=n-Q_i-(s_i-\sigma(i))<0,
$
so
$
s_i-\sigma(i)>n-Q_i.
$
Therefore
$
\binom{n-Q_i}{s_i-\sigma(i)}=0,
$
and hence $b_{i,\sigma(i)}=0$, 
so both sides of \eqref{E:11} vanish.
Therefore
\begin{align}\label{E:12}
[Y^E]H_p(X,Y)
&=X^m\sum_{\sigma\in S'_g}\sgn(\sigma)\prod_{i=1}^{g}b_{i,\sigma(i)}
\\
&=X^m\left(
\prod_{i=1}^{g}\binom{n}{Q_i}
\right)
\sum_{\sigma\in S'_g}
\sgn(\sigma)
\prod_{i=1}^{g}
\binom{n-Q_i}{s_i-\sigma(i)} \nonumber\\
&=X^m u_p \sum_{\sigma\in S'_g}\sgn(\sigma)\prod_{i=1}^{g}
\binom{n-Q_i}{s_i-\sigma(i)}. \nonumber
\end{align}

If $\sigma\in S_g\setminus S'_g$, then $\sigma(i)>s_i$ for some $i$, so
$
s_i-\sigma(i)<0.
$
By our convention for binomial coefficients,
$
\binom{n-Q_i}{s_i-\sigma(i)}=0.
$
We may therefore extend the sum in \eqref{E:12} from $S'_g$ to all of $S_g$. Using the Leibniz expansion of the determinant $\Delta_p$, we obtain
$$
\begin{aligned}
\sum_{\sigma\in S'_g}
\operatorname{sgn}(\sigma)
\prod_{i=1}^{g}
\binom{n-Q_i}{s_i-\sigma(i)}
&=
\sum_{\sigma\in S_g}
\operatorname{sgn}(\sigma)
\prod_{i=1}^{g}
\binom{n-Q_i}{s_i-\sigma(i)}=
\Delta_p.
\end{aligned}
$$
Combining this with \eqref{E:12}, we have 
$$
[Y^E]H_p(X,Y)=X^m u_p\Delta_p.
$$
{\bf (2).} Now let $R$ be a commutative ring and $a\in R$.
Applying the evaluation homomorphism  
$
\Z[X,Y]\longrightarrow R[Y]$ defined by
$X\longmapsto a,$
gives the specialization $H_p(a,Y)$.
Specialization cannot increase the $Y$-degree, so
$
\deg_Y H_p(a,Y)\leq E,
$
and the coefficient identity specializes to
$$
[Y^E]H_p(a,Y)=a^m u_p\Delta_p.
$$
In particular, when $m=0$, the polynomial $X^m=X^0=1$, 
and hence 
$$[Y^E]H_p(0,Y)=u_p\Delta_p.$$

\noindent 
{\bf (3).} Suppose now that $E<0$. For every $\sigma\in S_g$, 
$$
\sum_{i=1}^{g}\eta_{i,\sigma(i)}
=E-(d-1)e(\sigma)<0.
$$
If $\prod_{i=1}^g M_{i,\sigma(i)}(X,Y)\ne 0$, then 
$\eta_{i,\sigma(i)}\geq0$ for every $i$. 
Their sum would then be nonnegative, contradicting the preceding inequality. Thus every term in the Leibniz expansion of $H_p(X,Y)$ vanishes, and hence
$
H_p(X,Y)=0.
$

It remains to prove $\Delta_p=0$.
Consider a permutation $\sigma\in S_g$.
If $\sigma\not\in S'_g$, then 
$\sigma(i)>s_i$ for some $i$, so the corresponding Leibniz term 
in $\Delta_p$ is zero.
Now suppose $\sigma\in S'_g$.
If the corresponding Leibniz term were nonzero, then for every $i$, 
$\binom{n-Q_i}{s_i-\sigma(i)}\ne 0$.
Since $s_i-\sigma(i)\ge 0$, this implies 
$s_i-\sigma(i)\le n-Q_i,$
and therefore
$\eta_{i,\sigma(i)}=n-Q_i-s_i+\sigma(i)\ge 0.$
Since $\sigma\in S'_g$, Lemma \ref{L:E} gives 
$
E=\sum_{i=1}^{g}\eta_{i,\sigma(i)}\geq0,
$
contrary to the assumption $E<0$. Thus every term in the Leibniz
expansion of $\Delta_p$ is zero, and therefore
$
\Delta_p=0.
$
This proves part (3).
\end{proof}

\end{document}
